%
% backfill-run-lifecycle.spec.tex
%
% Backfill run lifecycle: a Z specification of one backfill run, from the
% singleton guard admitting it to the report it leaves behind, with cancel,
% daemon shutdown, a full queue, a hung worker, and a hung scan thread
% all in play (quarry-70zm).
%
% Check:        z-spec check docs/design/backfill-run-lifecycle.spec.tex
% Model-check:  probcli with -p MAX_INITIALISATIONS 100 (the command lines
%               are in backfill-run-lifecycle.spec.report.json)
%
% Must be .tex, not .md: the z-spec tools glob for .tex.
%

\documentclass[a4paper,10pt,fleqn]{article}
\usepackage[margin=1in]{geometry}
\usepackage{fuzz}
\usepackage{booktabs}

\newcommand{\field}[1]{\texttt{#1}}

\begin{document}

\title{Backfill Run Lifecycle: A Z Specification}
\author{Raymond H, for Punt Labs}
\date{October 2026}
\maketitle

% ===================================================================
\section{Scope and modelling choices}
\label{sec:scope}
% ===================================================================

One backfill run has four cooperating parts. The \emph{registry task}
(\field{task}) is what a client polls and what the singleton guard reads. The
\emph{scan thread} (\field{thr}) walks the transcripts in order. The
\emph{waiter} (\field{waiter}) is the coroutine that awaits the thread and
turns its result, or a cancellation, into the task's terminal state. The
\emph{worker} is the collection's single FIFO consumer; it runs the jobs the
thread submits, one transcript to a job, beside a \emph{foreign} job (a hook's
capture) that may occupy it.

The model has two transcripts and bounded counters: it checks the lifecycle,
not the arithmetic. A transcript is, in order, \emph{examined}, handed to the
queue as a \emph{job} (admitted, pending, running, finished), and
\emph{collected}, which gives it exactly one outcome.

\subsection{Two builds, one switch}

\field{fixed}~$=1$ is the lifecycle this design requires. \field{fixed}~$=0$ is
the lifecycle as it stood before: the second cancel escapes the wait, a
cancelled waiter leaves its job queued, a running job abandoned by the thread is
counted \emph{not attempted}, the guard reads the task alone, the cancel
overwrites the stop reason, and an unindexable transcript is recorded without
the racily-clean test. The same witnesses are checked against both; the
fixed build must reach none of them and the old build must reach each.

\subsection{Witnesses, not invariants}

A schema invariant that an operation can violate makes that operation
unsatisfiable, so ProB would report nothing rather than a violation. Each
property is therefore a latched witness in the set \field{wit}, set by the
transition that commits the bad step, and checked by asking ProB whether any
state with the witness in \field{wit} is reachable.

\subsection{Properties}

\begin{center}
\begin{tabular}{@{}cll@{}}
\toprule
 & Witness & Meaning \\
\midrule
L1 & \field{wOverlap} & the guard admitted a new run while the previous run's
  thread was still alive \\
L2 & \field{wLate} & a job wrote although its transcript was already counted
  as not attempted or not examined (the report did not count the write) \\
L3 & \field{wReason} & the task reports a reason other than the first stop
  reason \\
L4 & \field{wUncounted} & the run was reported while a transcript was still
  in flight, so it was counted in no outcome \\
L5 & \field{wRecord} & a transcript was recorded as unindexable although its
  mtime was racy or had changed \\
\bottomrule
\end{tabular}
\end{center}

L4 is also structural: \field{out} is a function from transcripts to outcomes,
so every transcript has exactly one, and the outcomes are disjoint.

\subsection{Non-goals}

\begin{itemize}
  \item Time. A timeout, the bound on the wait for the thread, and the
    admission wait are transitions that may fire whenever their guard holds;
    that they fire only after the configured seconds is not modelled.
  \item The transcript count, the chunk text, and the queue depth.
  \item The single strict read (L5): the read, the parse, and the verdict are
    one step here. That the code reads once is pinned by a test, not by this
    model.
  \item More than one foreign job, and more than one run at a time. The guard is
    modelled by a second start, \field{Start2}, that fires at most once.
\end{itemize}

% ===================================================================
\section{Types and constants}
\label{sec:types}
% ===================================================================

\begin{zed}
ZBOOL ::= ztrue | zfalse
\end{zed}

\begin{zed}
TASK ::= trunning | tcompleted | tfailed | tcancelled
\end{zed}

\begin{zed}
THR ::= talive | treturned
\end{zed}

The stop signal's reason. The first reason set is the one it keeps.

\begin{zed}
REASON ::= rnone | rcancel | rshutdown | rhung
\end{zed}

\begin{zed}
WAITER ::= wwaiting | wgone
\end{zed}

A transcript's job. \field{jadmit}: waiting for a queue slot, nothing enqueued.
\field{jpending}: queued. \field{jrunning}: the worker holds it.
\field{jstored}, \field{jnothing} (stored nothing: the text chunks to nothing),
\field{jfailed}: the worker finished it. \field{jskipped}: it will never write.

\begin{zed}
JOB ::= jnone | jadmit | jpending | jrunning
      | jstored | jnothing | jfailed | jskipped
\end{zed}

\begin{zed}
OUT ::= onone | oingested | oskipped | ofailed | onotattempted | onotexamined
\end{zed}

A transcript's content: ordinary; chunks to nothing with a stable mtime; chunks
to nothing with an mtime inside the racy window.

\begin{zed}
KIND ::= knormal | kstable | kracy
\end{zed}

\begin{zed}
WIT ::= wOverlap | wLate | wReason | wUncounted | wRecord
\end{zed}

\begin{axdef}
fixed : \nat
\where
fixed = 1
\end{axdef}

How the thread abandons a wait on a job, by the job's state: what the job
becomes and what the transcript is counted as. A queued job is cancelled so the
worker skips it; a running job continues and is counted \emph{failed} (it may
still land); the old build left the queued job in place and counted the running
one \emph{not attempted}.

\begin{axdef}
abortJob : JOB \pfun JOB \\
abortOut : JOB \pfun OUT
\where
\dom abortJob = \{ jadmit, jpending, jrunning \} \\
\dom abortOut = \{ jadmit, jpending, jrunning \} \\
abortJob~jadmit = jskipped \\
abortJob~jrunning = jrunning \\
abortOut~jadmit = onotattempted \\
abortOut~jpending = onotattempted \\
fixed = 1 \implies abortJob~jpending = jskipped \\
fixed = 0 \implies abortJob~jpending = jpending \\
fixed = 1 \implies abortOut~jrunning = ofailed \\
fixed = 0 \implies abortOut~jrunning = onotattempted
\end{axdef}

% ===================================================================
\section{State}
\label{sec:state}
% ===================================================================

\begin{schema}{Run}
task : TASK \\
rep, stop : REASON \\
thr : THR \\
waiter : WAITER \\
hung, gate, second, fbusy : ZBOOL \\
cancels : 0 \upto 2 \\
cur : 1 \upto 3 \\
job : 1 \upto 2 \fun JOB \\
out : 1 \upto 2 \fun OUT \\
kind : 1 \upto 2 \fun KIND \\
touched, unind : 1 \upto 2 \fun ZBOOL \\
wit : \power WIT
\end{schema}

\begin{itemize}
  \item \field{task}, \field{rep}: the registry task's status, and the stop
    reason the task state carries once the run is reported.
  \item \field{stop}: the first reason the stop signal was given.
  \item \field{thr}: the scan thread. \field{hung}: it is blocked in a step it
    controls (a read, a scrub, a directory listing) and takes no step until it
    returns from it.
  \item \field{waiter}: the coroutine that reports the run; \field{wgone} once
    the run is reported.
  \item \field{gate}: the queue's gate is closed (shutdown). \field{second}: the
    guard has admitted a second run. \field{fbusy}: a foreign job holds the
    worker.
  \item \field{cancels}: how many times the task was cancelled.
  \item \field{cur}: the transcript the thread is on; $3$ means the scan is
    complete.
  \item \field{touched}: the transcript's mtime changed after the snapshot.
    \field{unind}: the transcript is recorded as unindexable.
\end{itemize}

\begin{schema}{Init}
Run'
\where
task' = trunning \\
rep' = rnone \\
stop' = rnone \\
thr' = talive \\
waiter' = wwaiting \\
hung' = zfalse \\
gate' = zfalse \\
second' = zfalse \\
cancels' = 0 \\
cur' = 1 \\
\forall i : 1 \upto 2 @ job'~i = jnone \land out'~i = onone \land
  touched'~i = zfalse \land unind'~i = zfalse \\
wit' = \emptyset
\end{schema}

\field{kind} and \field{fbusy} are free in the initial state.

% ===================================================================
\section{Stop, cancel, and the foreign job}
\label{sec:stop}
% ===================================================================

Every operation below spells out what it leaves unchanged. (A shared frame
schema would itself be an operation to ProB, free to change the variables it
does not mention.) A stop reason is recorded only if none is set: the first
reason wins.

\begin{schema}{Cancel}
\Delta Run
\where
task = trunning \\
cancels < 2 \\
fixed = 1 \lor cancels = 0 \lor thr = treturned \lor waiter = wgone \\
cancels' = cancels + 1 \\
stop = rnone \implies stop' = rcancel \\
stop \neq rnone \implies stop' = stop \\
gate' = gate \\
task' = task \\
rep' = rep \\
thr' = thr \\
hung' = hung \\
cur' = cur \\
waiter' = waiter \\
job' = job \\
out' = out \\
touched' = touched \\
unind' = unind \\
kind' = kind \\
fbusy' = fbusy \\
second' = second \\
wit' = wit
\end{schema}

\begin{schema}{GateClose}
\Delta Run
\where
task = trunning \\
gate = zfalse \\
gate' = ztrue \\
cancels' = cancels \\
stop = rnone \implies stop' = rshutdown \\
stop \neq rnone \implies stop' = stop \\
task' = task \\
rep' = rep \\
thr' = thr \\
hung' = hung \\
cur' = cur \\
waiter' = waiter \\
job' = job \\
out' = out \\
touched' = touched \\
unind' = unind \\
kind' = kind \\
fbusy' = fbusy \\
second' = second \\
wit' = wit
\end{schema}


The old build's second cancel arrived while the waiter was awaiting the thread;
it escaped the wait, the task became cancelled at once, and whatever the thread
was in the middle of was left out of the report.

\begin{schema}{CancelAbandon}
\Delta Run
\where
fixed = 0 \\
task = trunning \\
cancels = 1 \\
waiter = wwaiting \\
thr = talive \\
cancels' = 2 \\
task' = tcancelled \\
rep' = rcancel \\
waiter' = wgone \\
out' = out \oplus \{ j : 1 \upto 2 | out~j = onone \land job~j = jnone @
  j \mapsto onotexamined \} \\
wit' = wit \cup \{ w : WIT |
  (w = wUncounted \land (\exists j : 1 \upto 2 @ out~j = onone \land job~j \neq jnone))
  \lor (w = wReason \land stop \neq rnone \land stop \neq rcancel) \} \\
stop' = stop \\
gate' = gate \\
thr' = thr \\
hung' = hung \\
cur' = cur \\
job' = job \\
touched' = touched \\
unind' = unind \\
kind' = kind \\
fbusy' = fbusy \\
second' = second
\end{schema}

\begin{schema}{FoeStart}
\Delta Run
\where
fbusy = zfalse \\
fbusy' = ztrue \\
kind' = kind \\
second' = second \\
task' = task \\
rep' = rep \\
stop' = stop \\
cancels' = cancels \\
gate' = gate \\
thr' = thr \\
hung' = hung \\
cur' = cur \\
waiter' = waiter \\
job' = job \\
out' = out \\
touched' = touched \\
unind' = unind \\
wit' = wit
\end{schema}

\begin{schema}{FoeEnd}
\Delta Run
\where
fbusy = ztrue \\
fbusy' = zfalse \\
kind' = kind \\
second' = second \\
task' = task \\
rep' = rep \\
stop' = stop \\
cancels' = cancels \\
gate' = gate \\
thr' = thr \\
hung' = hung \\
cur' = cur \\
waiter' = waiter \\
job' = job \\
out' = out \\
touched' = touched \\
unind' = unind \\
wit' = wit
\end{schema}

% ===================================================================
\section{The scan thread}
\label{sec:thread}
% ===================================================================

The thread checks the stop signal before every step it controls: before it
examines a transcript, before it reads one, before the scrub, and before the
capture file is written. \field{Examine} and \field{SkipInline} are those steps;
neither fires once a reason is set. A thread that is blocked inside a step
(\field{Hang}) takes no step until \field{Unhang}.

\begin{schema}{Examine}
\Delta Run
\where
thr = talive \\
hung = zfalse \\
cur \in 1 \upto 2 \\
stop = rnone \\
job~cur = jnone \\
out~cur = onone \\
job' = job \oplus \{ cur \mapsto jadmit \} \\
task' = task \\
rep' = rep \\
stop' = stop \\
cancels' = cancels \\
gate' = gate \\
thr' = thr \\
hung' = hung \\
cur' = cur \\
waiter' = waiter \\
out' = out \\
touched' = touched \\
unind' = unind \\
kind' = kind \\
fbusy' = fbusy \\
second' = second \\
wit' = wit
\end{schema}

\begin{schema}{SkipInline}
\Delta Run
\where
thr = talive \\
hung = zfalse \\
cur \in 1 \upto 2 \\
stop = rnone \\
job~cur = jnone \\
out~cur = onone \\
out' = out \oplus \{ cur \mapsto oskipped \} \\
cur' = cur + 1 \\
task' = task \\
rep' = rep \\
stop' = stop \\
cancels' = cancels \\
gate' = gate \\
thr' = thr \\
hung' = hung \\
waiter' = waiter \\
job' = job \\
touched' = touched \\
unind' = unind \\
kind' = kind \\
fbusy' = fbusy \\
second' = second \\
wit' = wit
\end{schema}

\begin{schema}{Hang}
\Delta Run
\where
thr = talive \\
hung = zfalse \\
cur \in 1 \upto 2 \\
job~cur = jnone \\
out~cur = onone \\
hung' = ztrue \\
thr' = thr \\
task' = task \\
rep' = rep \\
stop' = stop \\
cancels' = cancels \\
gate' = gate \\
cur' = cur \\
waiter' = waiter \\
job' = job \\
out' = out \\
touched' = touched \\
unind' = unind \\
kind' = kind \\
fbusy' = fbusy \\
second' = second \\
wit' = wit
\end{schema}

\begin{schema}{Unhang}
\Delta Run
\where
hung = ztrue \\
hung' = zfalse \\
thr' = thr \\
task' = task \\
rep' = rep \\
stop' = stop \\
cancels' = cancels \\
gate' = gate \\
cur' = cur \\
waiter' = waiter \\
job' = job \\
out' = out \\
touched' = touched \\
unind' = unind \\
kind' = kind \\
fbusy' = fbusy \\
second' = second \\
wit' = wit
\end{schema}

The thread returns when the scan is complete, or when it sees a stop reason
between transcripts; every transcript it did not reach is counted \emph{not
examined}.

\begin{schema}{ReturnDone}
\Delta Run
\where
thr = talive \\
hung = zfalse \\
cur = 3 \\
thr' = treturned \\
hung' = hung \\
task' = task \\
rep' = rep \\
stop' = stop \\
cancels' = cancels \\
gate' = gate \\
cur' = cur \\
waiter' = waiter \\
job' = job \\
out' = out \\
touched' = touched \\
unind' = unind \\
kind' = kind \\
fbusy' = fbusy \\
second' = second \\
wit' = wit
\end{schema}

\begin{schema}{ReturnStopped}
\Delta Run
\where
thr = talive \\
hung = zfalse \\
cur \in 1 \upto 2 \\
stop \neq rnone \\
job~cur = jnone \\
thr' = treturned \\
hung' = hung \\
out' = out \oplus \{ j : 1 \upto 2 | out~j = onone @ j \mapsto onotexamined \} \\
task' = task \\
rep' = rep \\
stop' = stop \\
cancels' = cancels \\
gate' = gate \\
cur' = cur \\
waiter' = waiter \\
job' = job \\
touched' = touched \\
unind' = unind \\
kind' = kind \\
fbusy' = fbusy \\
second' = second \\
wit' = wit
\end{schema}

% ===================================================================
\section{The queue and the worker}
\label{sec:queue}
% ===================================================================

\field{Admit} takes a slot (a closed gate admits nothing). \field{AdmitTimeout}
ends the admission wait with the transcript not attempted; nothing was
enqueued. \field{AbortWait} is the thread leaving a wait because a stop reason
was set, honored within a poll. \field{GiveUp} is the finish timeout: it sets
the reason \field{rhung} and leaves the wait the same way.

\begin{schema}{Admit}
\Delta Run \\ i? : 1 \upto 2
\where
job~i? = jadmit \\
gate = zfalse \\
job' = job \oplus \{ i? \mapsto jpending \} \\
task' = task \\
rep' = rep \\
stop' = stop \\
cancels' = cancels \\
gate' = gate \\
thr' = thr \\
hung' = hung \\
cur' = cur \\
waiter' = waiter \\
out' = out \\
touched' = touched \\
unind' = unind \\
kind' = kind \\
fbusy' = fbusy \\
second' = second \\
wit' = wit
\end{schema}

\begin{schema}{AdmitTimeout}
\Delta Run
\where
thr = talive \\
hung = zfalse \\
cur \in 1 \upto 2 \\
stop = rnone \\
job~cur = jadmit \\
job' = job \oplus \{ cur \mapsto jskipped \} \\
out' = out \oplus \{ cur \mapsto onotattempted \} \\
cur' = cur + 1 \\
task' = task \\
rep' = rep \\
stop' = stop \\
cancels' = cancels \\
gate' = gate \\
thr' = thr \\
hung' = hung \\
waiter' = waiter \\
touched' = touched \\
unind' = unind \\
kind' = kind \\
fbusy' = fbusy \\
second' = second \\
wit' = wit
\end{schema}

\begin{schema}{AbortWait}
\Delta Run
\where
thr = talive \\
hung = zfalse \\
cur \in 1 \upto 2 \\
stop \neq rnone \\
job~cur \in \dom abortJob \\
job' = job \oplus \{ cur \mapsto abortJob~(job~cur) \} \\
out' = out \oplus \{ cur \mapsto abortOut~(job~cur) \} \\
cur' = cur + 1 \\
task' = task \\
rep' = rep \\
stop' = stop \\
cancels' = cancels \\
gate' = gate \\
thr' = thr \\
hung' = hung \\
waiter' = waiter \\
touched' = touched \\
unind' = unind \\
kind' = kind \\
fbusy' = fbusy \\
second' = second \\
wit' = wit
\end{schema}

\begin{schema}{GiveUp}
\Delta Run
\where
thr = talive \\
hung = zfalse \\
cur \in 1 \upto 2 \\
stop = rnone \\
job~cur \in \{ jpending, jrunning \} \\
stop' = rhung \\
cancels' = cancels \\
gate' = gate \\
job' = job \oplus \{ cur \mapsto abortJob~(job~cur) \} \\
out' = out \oplus \{ cur \mapsto abortOut~(job~cur) \} \\
cur' = cur + 1 \\
task' = task \\
rep' = rep \\
thr' = thr \\
hung' = hung \\
waiter' = waiter \\
touched' = touched \\
unind' = unind \\
kind' = kind \\
fbusy' = fbusy \\
second' = second \\
wit' = wit
\end{schema}

The worker takes the oldest pending job when no foreign job holds it. It runs a
job to one of three ends. A job that stores writes; a job whose transcript was
already counted \emph{not attempted} or \emph{not examined} that then writes is
the late write L2 forbids.

\begin{schema}{WorkerStart}
\Delta Run \\ i? : 1 \upto 2
\where
job~i? = jpending \\
fbusy = zfalse \\
\forall j : 1 \upto 2 @ job~j \neq jrunning \\
\forall j : 1 \upto 2 | j < i? @ job~j \neq jpending \\
job' = job \oplus \{ i? \mapsto jrunning \} \\
task' = task \\
rep' = rep \\
stop' = stop \\
cancels' = cancels \\
gate' = gate \\
thr' = thr \\
hung' = hung \\
cur' = cur \\
waiter' = waiter \\
out' = out \\
touched' = touched \\
unind' = unind \\
kind' = kind \\
fbusy' = fbusy \\
second' = second \\
wit' = wit
\end{schema}

\begin{schema}{WorkerStores}
\Delta Run \\ i? : 1 \upto 2
\where
job~i? = jrunning \\
kind~i? = knormal \\
job' = job \oplus \{ i? \mapsto jstored \} \\
wit' = wit \cup \{ w : WIT | w = wLate \land out~i? \in \{ onotattempted, onotexamined \} \} \\
task' = task \\
rep' = rep \\
stop' = stop \\
cancels' = cancels \\
gate' = gate \\
thr' = thr \\
hung' = hung \\
cur' = cur \\
waiter' = waiter \\
out' = out \\
touched' = touched \\
unind' = unind \\
kind' = kind \\
fbusy' = fbusy \\
second' = second
\end{schema}

\begin{schema}{WorkerStoresNothing}
\Delta Run \\ i? : 1 \upto 2
\where
job~i? = jrunning \\
kind~i? \neq knormal \\
job' = job \oplus \{ i? \mapsto jnothing \} \\
task' = task \\
rep' = rep \\
stop' = stop \\
cancels' = cancels \\
gate' = gate \\
thr' = thr \\
hung' = hung \\
cur' = cur \\
waiter' = waiter \\
out' = out \\
touched' = touched \\
unind' = unind \\
kind' = kind \\
fbusy' = fbusy \\
second' = second \\
wit' = wit
\end{schema}

\begin{schema}{WorkerFails}
\Delta Run \\ i? : 1 \upto 2
\where
job~i? = jrunning \\
job' = job \oplus \{ i? \mapsto jfailed \} \\
task' = task \\
rep' = rep \\
stop' = stop \\
cancels' = cancels \\
gate' = gate \\
thr' = thr \\
hung' = hung \\
cur' = cur \\
waiter' = waiter \\
out' = out \\
touched' = touched \\
unind' = unind \\
kind' = kind \\
fbusy' = fbusy \\
second' = second \\
wit' = wit
\end{schema}

% ===================================================================
\section{Collecting a transcript}
\label{sec:collect}
% ===================================================================

The thread collects a finished job into the transcript's one outcome and moves
on. A transcript whose job stored nothing is recorded as unindexable only if
its mtime is unchanged since the snapshot and, in the fixed build, outside the
racy window as well (L5). \field{Touch} is the file changing under the scan.

\begin{schema}{CollectStored}
\Delta Run
\where
thr = talive \\
hung = zfalse \\
cur \in 1 \upto 2 \\
out~cur = onone \\
job~cur = jstored \\
out' = out \oplus \{ cur \mapsto oingested \} \\
cur' = cur + 1 \\
task' = task \\
rep' = rep \\
stop' = stop \\
cancels' = cancels \\
gate' = gate \\
thr' = thr \\
hung' = hung \\
waiter' = waiter \\
job' = job \\
touched' = touched \\
unind' = unind \\
kind' = kind \\
fbusy' = fbusy \\
second' = second \\
wit' = wit
\end{schema}

\begin{schema}{CollectFailed}
\Delta Run
\where
thr = talive \\
hung = zfalse \\
cur \in 1 \upto 2 \\
out~cur = onone \\
job~cur = jfailed \\
out' = out \oplus \{ cur \mapsto ofailed \} \\
cur' = cur + 1 \\
task' = task \\
rep' = rep \\
stop' = stop \\
cancels' = cancels \\
gate' = gate \\
thr' = thr \\
hung' = hung \\
waiter' = waiter \\
job' = job \\
touched' = touched \\
unind' = unind \\
kind' = kind \\
fbusy' = fbusy \\
second' = second \\
wit' = wit
\end{schema}

\begin{schema}{CollectRecord}
\Delta Run
\where
thr = talive \\
hung = zfalse \\
cur \in 1 \upto 2 \\
out~cur = onone \\
job~cur = jnothing \\
touched~cur = zfalse \\
kind~cur = kstable \lor (fixed = 0 \land kind~cur = kracy) \\
out' = out \oplus \{ cur \mapsto oskipped \} \\
unind' = unind \oplus \{ cur \mapsto ztrue \} \\
touched' = touched \\
cur' = cur + 1 \\
wit' = wit \cup \{ w : WIT | w = wRecord \land kind~cur = kracy \} \\
task' = task \\
rep' = rep \\
stop' = stop \\
cancels' = cancels \\
gate' = gate \\
thr' = thr \\
hung' = hung \\
waiter' = waiter \\
job' = job \\
kind' = kind \\
fbusy' = fbusy \\
second' = second
\end{schema}

\begin{schema}{CollectNoRecord}
\Delta Run
\where
thr = talive \\
hung = zfalse \\
cur \in 1 \upto 2 \\
out~cur = onone \\
job~cur = jnothing \\
touched~cur = ztrue \lor kind~cur = knormal \lor
  (fixed = 1 \land kind~cur = kracy) \\
out' = out \oplus \{ cur \mapsto oskipped \} \\
cur' = cur + 1 \\
task' = task \\
rep' = rep \\
stop' = stop \\
cancels' = cancels \\
gate' = gate \\
thr' = thr \\
hung' = hung \\
waiter' = waiter \\
job' = job \\
touched' = touched \\
unind' = unind \\
kind' = kind \\
fbusy' = fbusy \\
second' = second \\
wit' = wit
\end{schema}

\begin{schema}{Touch}
\Delta Run \\ i? : 1 \upto 2
\where
out~i? = onone \\
touched~i? = zfalse \\
touched' = touched \oplus \{ i? \mapsto ztrue \} \\
unind' = unind \\
task' = task \\
rep' = rep \\
stop' = stop \\
cancels' = cancels \\
gate' = gate \\
thr' = thr \\
hung' = hung \\
cur' = cur \\
waiter' = waiter \\
job' = job \\
out' = out \\
kind' = kind \\
fbusy' = fbusy \\
second' = second \\
wit' = wit
\end{schema}

% ===================================================================
\section{Reporting the run, and the guard}
\label{sec:report}
% ===================================================================

\field{Settle} is the waiter reporting a run whose thread has returned. A
cancelled run is \field{tcancelled}; otherwise it completes or fails. The task
carries the first stop reason (L3); the old build carried \field{rcancel}
whenever the task was cancelled.

\begin{schema}{Settle}
\Delta Run
\where
task = trunning \\
waiter = wwaiting \\
thr = treturned \\
cancels = 0 \implies task' \in \{ tcompleted, tfailed \} \\
cancels > 0 \implies task' = tcancelled \\
(fixed = 1 \lor cancels = 0) \implies rep' = stop \\
(fixed = 0 \land cancels > 0) \implies rep' = rcancel \\
waiter' = wgone \\
wit' = wit \cup \{ w : WIT |
  (w = wReason \land stop \neq rnone \land rep' \neq stop)
  \lor (w = wUncounted \land (\exists j : 1 \upto 2 @ out~j = onone)) \} \\
stop' = stop \\
cancels' = cancels \\
gate' = gate \\
thr' = thr \\
hung' = hung \\
cur' = cur \\
job' = job \\
out' = out \\
touched' = touched \\
unind' = unind \\
kind' = kind \\
fbusy' = fbusy \\
second' = second
\end{schema}

\field{Bound} is the bounded wait: the thread is blocked in a step the stop
signal cannot interrupt, the stop reason has been set, and the waiter stops
waiting. The task is reported at once, the thread is still alive, and the guard
reads the thread as well as the task, so it refuses a second run until the
thread returns.

\begin{schema}{Bound}
\Delta Run
\where
fixed = 1 \\
task = trunning \\
waiter = wwaiting \\
thr = talive \\
hung = ztrue \\
stop \neq rnone \\
cancels = 0 \implies task' = tfailed \\
cancels > 0 \implies task' = tcancelled \\
rep' = stop \\
waiter' = wgone \\
out' = out \oplus \{ j : 1 \upto 2 | out~j = onone @ j \mapsto onotexamined \} \\
wit' = wit \cup \{ w : WIT |
  w = wUncounted \land (\exists j : 1 \upto 2 @ out~j = onone \land job~j \neq jnone) \} \\
stop' = stop \\
cancels' = cancels \\
gate' = gate \\
thr' = thr \\
hung' = hung \\
cur' = cur \\
job' = job \\
touched' = touched \\
unind' = unind \\
kind' = kind \\
fbusy' = fbusy \\
second' = second
\end{schema}

\field{Start2} is a second backfill request reaching the singleton guard once
the task is no longer running. The fixed guard also waits for the thread.

\begin{schema}{Start2}
\Delta Run
\where
second = zfalse \\
task \neq trunning \\
fixed = 0 \lor thr = treturned \\
second' = ztrue \\
kind' = kind \\
fbusy' = fbusy \\
wit' = wit \cup \{ w : WIT | w = wOverlap \land thr = talive \} \\
task' = task \\
rep' = rep \\
stop' = stop \\
cancels' = cancels \\
gate' = gate \\
thr' = thr \\
hung' = hung \\
cur' = cur \\
waiter' = waiter \\
job' = job \\
out' = out \\
touched' = touched \\
unind' = unind
\end{schema}

% ===================================================================
\section{Properties checked}
\label{sec:props}
% ===================================================================

For \field{fixed}~$=1$, ProB finds no reachable state with any witness in
\field{wit}, no deadlock before the run is reported, and from every reachable
state a state in which the task is terminal, the thread has returned and
every transcript has an outcome is reachable (\field{AG(EF(done))}). For
\field{fixed}~$=0$, each witness is reachable; the reports record the shortest
path ProB found for each.

\end{document}
